\section{The Maven pipeline}
\label{sec:C.1}
This is the translation \S\ref{sec:18.3} runs on the Maven Smart System: for each article, the terms derived for the pipeline, and the terms considered and rejected and why. The articles that yield none are listed at the end, and the Geneva Conventions and Additional Protocol I are translated article by article in \S\ref{sec:C.2}.

\runin{Table C.1. What the instruments yield for the Maven pipeline}

\textbf{UN Charter Art. 2(4); exceptions Art. 51 and Chapter VII.} Don't take part in the use of armed force, except for corroborated defense, authorized force, the genocide exception or the occupation exception; and, where the belligerent's constitution requires legislative authorization, not without it. A disinterested finding for defense means a judgment of the International Court of Justice or of a regional court with adversarial process, a commission of inquiry mandated by a body no party to the conflict controls, a Security Council determination, or a two-thirds General Assembly resolution where the facts aren't in dispute. For the genocide and occupation exceptions, a merits judgment or opinion of the Court, or two findings of different kinds, one of them from a body that took evidence; a provisional-measures order counts as one, for genocide only. Rejected: lift on the state's own Art. 51 claim, because the Charter allows it until the Council acts; this floor doesn't, and says so. Rejected: require Council authorization for all force, because a veto is a stake, and the rule would forbid Ukraine's defense. Rejected: ask the machine to rule on constitutionality, because courts won't, and it shouldn't; the clause reads a marker.

\textbf{Art. 3 (life, liberty, security of person), read with AP I Arts. 48 and 52 (the machine-selection departure).} Don't resolve a description of a person into a location, unless it identifies the person by conduct in hostilities, with disinterested corroboration, the case falls within the review ratio, and the location isn't a residence. Don't time an action to presence at a residence. Rejected: don't generate targets at all, because this is now the default under the first entry, lifted only by its exceptions. Rejected for one case: bar all machine selection of persons for force, because a present, formed, recorded and halt-able system stopping an attack under way can answer for the selection where a munition can't.

\textbf{Art. 3 (life), read with AP I Art. 54 and Rome Statute Art. 8(2)(b)(xxv).} Don't take part in a siege, blockade or attack on objects indispensable to survival whose effect is to deprive a civilian population of food, water or medical care. Rejected: don't take part in any blockade, because a naval blockade of military supply is lawful under the law of armed conflict. The term reaches the effect on the population, not the form.

\textbf{AP I Arts. 51(4)–(5), 57(2)(a)(iii) (indiscriminate attack; proportionality).} Operator term: the action threshold is a published minimum under the commander's proportionality judgment, and the deployment states it. Model term: don't resolve a location for an attack the deployment's own threshold finds excessive. Rejected: ask the machine to make the proportionality judgment, because the law leaves it to a reasonable commander, and the threshold sits under that judgment, not in place of it.

\textbf{GC IV Arts. 18–19; AP I Arts. 12, 61–67 (medical units; civil defence).} Don't resolve a location at a medical unit or a civil-defence organization, its personnel, buildings or shelters, or time an action to presence at one, absent a disinterested finding that it is used for acts harmful to the enemy. Rejected: lift on the operator's claim of misuse, because that claim is the operator's own, and Art. 19 GC IV and Art. 65 AP I require a warning and a time limit before protection ceases.

\textbf{AP I Art. 53 (cultural objects and places of worship).} Don't resolve a location at a historic monument, work of art or place of worship that constitutes the cultural or spiritual heritage of peoples, or time an action to presence at one, absent a disinterested finding that it is used in support of the military effort. Rejected: lift on the operator's claim of military use, because that claim is the operator's own.

\textbf{AP I Arts. 59–60 (non-defended localities; demilitarized zones).} Don't resolve a location in a locality declared non-defended or a zone agreed demilitarized, unless a disinterested body finds the article's conditions no longer met. Rejected: lift on the operator's claim that the conditions lapsed, because that claim is the operator's own.

\textbf{AP I Art. 56 (works containing dangerous forces).} Don't resolve a location at a dam, dyke or nuclear generating station. Rejected: none, because the article's own exceptions require a determination the pipeline can't make.

\textbf{GC IV Art. 33; AP I Art. 51(2), 51(6) (collective punishment; reprisals against civilians; terror).} Don't take part in an action whose stated purpose is to punish a population for the acts of others or to spread terror among it. Rejected: a term over the effect, terror caused, because every bombing terrifies, and the article's line is primary purpose. A model can find purpose only where it is stated, so the term reads the stated purpose and the action threshold carries the rest.

\textbf{Art. 3 (life), read with the rule on machine estimates below and AP I Art. 48 (the machine-selection departure).} Don't rank people by estimated affiliation. Rejected: derive the term from Art. 11(1) by analogy, because the presumption of innocence governs punishment, targeting isn't punishment, and the law of armed conflict permits status-based targeting of organized armed group members; the ground is a claim about machines, stated below, not a translation of an article. Rejected: don't score anything, because scoring a site or vehicle for military function is not a finding about a person.

\textbf{Art. 12 (no arbitrary interference with home or family).} Shares “don't time an action to presence at a residence” with Art. 3. Rejected: don't hold location data about homes, because the no-strike list needs exactly that data.

\textbf{Arts. 8 and 10 (effective remedy; fair determination), by analogy.} Operator term: don't operate above the review ratio. Rejected: a model-side term requiring review, because the model can't establish whether review occurred.

\textbf{Arts. 2 and 7 (non-discrimination; equal protection).} Don't use ethnicity, religion or national origin as inputs to a person's score. Rejected: don't record those attributes at all, because auditing a pipeline for discriminatory effect needs them.

\textbf{Art. 19 (opinion and expression).} Don't use a person's opinions, published or private, as evidence of affiliation. Rejected: don't process a person's communications at all, because an intercepted operational order is evidence of a role in hostilities, which is conduct.

\textbf{Arts. 4–6, 9, 13–18, 20–21; GC I–III, all but Common Art. 3; GC IV Arts. 1–2, 4–17, 20–22, 24–32, 35–48, 50–52, 54, 56–58, 60–146, 148–159; AP I Arts. 1(1)–(3), 2–11, 13–34, 36, 38–47, 49–50, 55, 68–84, 86–102.} None: no act this pipeline performs engages them. Art. 17's property clause is left to IHL's distinction rule; Art. 9 becomes live in the removal case. For the Geneva articles: none for this deployment, with Appendix~\ref{sec:C} giving the article-by-article reasons. GC IV Arts. 23, 34, 49, 53, 55, 59 and 147, AP I Arts. 1(4), 35, 37, 58 and 85, and Common Art. 3 are written in full in Appendix~\ref{sec:C} and take their verdicts from it.
